% from main.tex
%%%%%%%%%%  scan idx 494  |  folio 481  %%%%%%%%%%
En
appliquant $i^*\circ j'^*$, on trouve la suite exacte :
\[
0\longrightarrow (p''\circ j'\circ i)^*(\Omega^1_{V'/Z})\longrightarrow (j'\circ i)^*(\Omega^1_{V''/Z})\longrightarrow (j'\circ i)^*(\Omega^1_{V''/V'})\longrightarrow0.
\]
Comme $j'^*(\Omega^1_{V''/V'})=\Omega^1_{V/Y}$ d'après l'hypothèse (G), on obtient :
\begin{equation}\tag{2.6.2}
0\longrightarrow f^*(j^*\Omega^1_{V'/Z})\longrightarrow (j'\circ i)^*(\Omega^1_{V''/Z})\longrightarrow i^*(\Omega^1_{V/Y})\longrightarrow0.
\end{equation}
En combinant (2.6.1) et (2.6.2), on est ramené à voir que le diagramme :
\[
\begin{tikzcd}[column sep=large,row sep=large]
0 \arrow[r] & f^*(J/J^2) \arrow[r] \arrow[d,"f^*(d)"'] & K/K^2 \arrow[r] \arrow[d,"d"] & I/I^2 \arrow[r] \arrow[d,"d"] & 0\\
0 \arrow[r] & f^*(j^*\Omega^1_{V'/Z}) \arrow[r] & (j'\circ i)^*\Omega^1_{V''/Z} \arrow[r] & i^*\Omega^1_{V/Y} \arrow[r] & 0
\end{tikzcd}
\]
est commutatif, ce qui se fait localement sans difficulté en revenant aux définitions des morphismes en jeu. De la suite exacte de complexes ainsi obtenue, on déduit le triangle distingué annoncé dans la catégorie dérivée.

\textbf{Corollaire 2.7.} Avec les notations et les hypothèses de 2.6, on a dans $K^{\bullet}(X)$ l'égalité suivante:
\[
T_{g\circ f}=T_f+f^*(T_g).
\]

\subsection*{3. Théorème de Riemann--Roch : énoncé}

\textbf{3.1.} Dans le numéro 3, $Y$ est un schéma quasi-compact ayant un faisceau ample.

Soit $f:X\to Y$ un morphisme projectif, d'intersection complète (1.1); on suppose de plus que la dimension relative virtuelle de $f$ (1.9) est
